L'església primitiva

L'any 1620, don Pau Sureda i Campfullós, senyor de l'alqueria de Sant Martí, obtingué permís per crear un nucli de població a les seves terres de Son Pere Jaume. Com que els seus habitants eren francs de les contribucions al Regne i altres imposts i només havien de pagar les taxes corresponents als senyors de Sant Martí, aquella vila va ser coneguda com Vilafranca de Sant Martí. Amb el temps es perdé l'afegitó "de Sant Martí" i l'any 1916, arrel d'una llei que feia obligatori afegir apel·latius als pobles d'Espanya que tenguéssin el mateix nom, se li imposà el "de Bonany". 

Aquell mateix 1620 s'inicià la construcció a la vila la primera església, d'estil renaixentista, que s'acabà l'any 1631; la primera missa s'hi digué el dia de santa Bàrbara (4 de desembre) d'aquell any. Es creu que ocupava el solar dels dos primers trams de l'església actual, començant pel portal principal. Va ser dedicada a santa Bàrbara i sant Martí, els sants patrons de la família Sureda, senyors de Sant Martí.

La seva construcció va sofrir una forta oposició dels petrencs, de quina parròquia depenia la vila, perquè a Petra estaven construint una església nova i volien evitar que es desviassin recursos cap a la de Vilafranca. La major part de la despesa de la construcció la pagà la família Sureda. Els vilafranquers, a banda de la mà d'obra, aportaren els diners de la venda dels animals que trobaven perduts pel camp, de la qual cosa els Jurats de Petra se n'exclamaven: ``lo cual mos pareix meravella i cosa nova may usada.'' 

Els primers anys venia un capellà de Petra a cel·lebrar els actes de culte, fins que l'any 1685 s'hi instituí una vicaria, dependent de la parròquia de Petra. El primer vicari que anomenà el rector de Petra va ser Miquel Riera, un nebot seu que encara no havia estat ordenat sacerdot, per la qual cosa, per exemple, els primers matrimonis que es cel·lebraren a Vilafranca els hagué d'oficiar un frare franciscà. El terreny destinat a cementiri era l'esplanada que envolta l'església pel sud i l'est, anomenat "El Sagrat" perquè era fins on s'estenia el dret d'assil de l'església. 

El dia de Sant Agustí, 28 d'agost, del 1685 es diposità a l'església per primer cop el Santíssim Sagrament, portat en processó des de Petra, i des de llavors es recordà aquesta efemèride cel·lebrant l'anomenada festa de la Reserva; aquesta va ser la festa grossa de l'estiu a la vila fins que, durant la segona meitat del segle XX, va ser substituïda poc a poc per les festes de la Beata, al voltant del 28 de juliol. 

Vilafranca no va ser parròquia independent fins el 1913.


L'església antiga

Amb el creixement de la població de Vilafranca, el temple es va anar quedant petit. A voltant del 1700 es començà la construcció d'una església més gran a partir de la primera. El nou temple va ser d'estil barroc, amb una nau única, bòveda de canó, capelles als dos costats de la nau i una rosassa a la façana principal. Tenia 35 m de llarg, 17 m d'ample i 16 m d'alçada, cinc capelles a cada costat  i dos portals: el principal, o ``de les dones'', i el lateral, o ``dels homes''. El buc de l'església s'acabà l'any 1738, i  la façana el 1741. 

Un dels principals impulsors de l'església nova va ser el canonge Domingo Sureda de Sant Martí, fill del senyor de Sant Martí i que arribà a rector de la Universitat Luliana i Vicari General de la diòcesi. Entre altres coses, donà a l'església de Vilafranca una relíquia de la Vera Creu que arrabassà de la que es conserva a la Seu. 

Al voltant del 1775 s'encarregà el retaule major i el sagrari monumental a l'escultor Mateu Colom, assessorat per D. Jeroni Bestard, amic del marquès de Sant Martí i artista de gran prestigi a Mallorca. Bestard volia que el retaule fos neoclàssic, però Colom el construí barroc, que encara era la moda. A Bestard no li agradà i arran d'una visita comentà: "tuve el desconsuelo de verlo alborotado por las quiméricas ideas del artífice que lo construyó". El  retaule de Colom és la part central de l'actual retaule major.

Naturalment, com en tota casa, les obres continuaren més enllà de l'acabament de l'obra principal. Així, per exemple:

* l'any 1802 es col·locà la rosassa de vidre (que s'hagué de reconstruir poc després, perquè l'any 1804 va ser destrossada per un llamp);
* el 1803 es construí el cancell del portal lateral i el 1804 el del portal principal;
* entre el 1817 i el 1842 es construí el campanar: l'antic campanar eren les dues finestres tapiades que es veuen a l'angle superior esquerre de la façana principal;
* el 1907 es rebaixà el paviment uns 60 cm per fer el temple més esbelt;
* el mateix 1907 es substituí la pila baptismal antiga, del 1685, per una de nova, dissenyada pel marquès de Vivot  i llaurada per Antoni Rotger, i s'obrí la finestra del baptisteri i s'hi col·locà el  vitrall.

L'església nova

L'any 1935,  en el 250è anniversari de la Reserva, començaren les obres d'ampliació del temple que donaren lloc a l'actual. S'afegiren l'absis, els vitralls i el creuer amb una cúpula central i quatre cúpules més petites. La longitud final del temple quedà en els actuals 50 m; la cúpula central és de 30 m de diàmetre a la base i s'alça fins al 28 m sobre el trespol. Les obres acabaren l'any 1941. La direcció de l'obra va ser de l'arquitecte Gulliem Ferragut.

Sobre la cúpula del creuer, l'any 1942 s’hi instal·là una escultura del Cor de Jesús de 6 m d'alt, comptant la peanya, feta de ciment portland armat, obra de l’escultor Miquel Vadell ajudat per Joan Bauzà Font, el mestre d'obres local. A continuació s'amplià el retaule del presbiteri, que quedava petit a l'absis de l'església nova, afegint als tres cossos del retaule antic un cos més a cada costat. L'autor de l'ampliació va ser l'escultor Bartomeu Amorós. L'altar de pedra nou es va estrenar per la festa de la Reserva del 1947.



Rovissos

Un dels motius de la popularitat de la Mare de Déu del Carme és el seu paper com a patrona dels mariners i pescadors. Segons una llegenda, la Verge salvà un vaixell d’enfonsar-se enmig d’una tempesta quan un dels mariners llançà a la mar un escapulari carmelita que duia al coll. Aquest patronatge fou promogut al segle XVIII per l’almirall mallorquí Antoni Barceló, el “capità Toni”, que exigia als seus mariners encomanar-se a la Mare de Déu del Carme i duia la seva imatge al buc insígnia de la seva armada.//

Els tres cossos centrals formen el retaule de l’església antiga, obra de Mateu Colom (s. XVIII). Quan a mitjan segle XX es traslladà a l’absis de la nova quedà petit i l’ampliàren amb un cos més a cada costat . Als portals de les sagristies JHS amb creu i tres claus, emblema que els jesuites tenen incorporat al seu escut, i a l’altre AM d’Ave Maria. Fotos?

